
% Exposé VII: back to the volume's Sym and Proj
\renewcommand{\Sym}{\operatorname{Sym}}
\renewcommand{\Projz}{\operatorname{Proj}_0}
\section*{Exposé VII\\[3pt] IMMERSIONS RÉGULIÈRES ET CALCUL DU \texorpdfstring{$K^{\bullet}$}{K0} D'UN SCHÉMA ÉCLATÉ}
\begin{center}
\uline{par Pierre BERTHELOT}
\end{center}

\section*{1. Généralités sur les immersions régulières}

La notion classique d'immersion régulière (EGA IV 16.9.2) s'avérant inadéquate dans le cas non noethérien, nous allons définir une notion plus faible, mais donnant lieu aux propriétés souhaitées, et qui équivaut à la définition usuelle dans le cas noethérien.

\textbf{Définition 1.1.} \uline{Soient} $A$ \uline{un anneau commutatif}, $E$ \uline{un} $A$-\uline{module projectif de type fini}, $u:E\longrightarrow A$ \uline{un homomorphisme} $A$-\uline{linéaire}. \uline{On dit que} $u$ \uline{est régulier si le complexe de Koszul} $K_\bullet(u)$ \uline{défini par} $u$ \uline{est acyclique en dimension} $>0$, \uline{donc est une résolution de} $B=A/u(E)$.

Rappelons que le complexe de Koszul associé à $u$ est défini comme suit: le $A$-module gradué $K_\bullet(u)$ est l'algèbre extérieure $\Lambda(E)$, et l'opérateur bord la multiplication intérieure par $u$ considéré comme élément du dual $\check E$ : c'est une antidérivation de degré $-1$ de $\Lambda(E)$, égale à $u$ sur $E$.

\textbf{Lemme 1.2.} \uline{Soit} $u:E\longrightarrow A$ \uline{un homomorphisme régulier}. \uline{Soient} $A'$ \uline{une} $A$-\uline{algèbre}, $u':E'\longrightarrow A'$ \uline{l'homomorphisme déduit de} $u$ \uline{par changement de base}; \uline{on suppose} $u'$ \uline{régulier}. \uline{On pose} $J=u(E)$, $J'=JA'=u(E')$. \uline{Alors on a les relations suivantes pour tout entier} $n\geq0$:
\begin{align}
 J^n/J^{n+1}\otimes_AA'&\xrightarrow{\sim}J'^n/J'^{n+1}, \tag{1.2.1}\\
 \Tor_i^A(J^n/J^{n+1},A')&=0\quad\text{\uline{pour} }i\geq1; \tag{1.2.2}\\
 J^n\otimes_AA'&\xrightarrow{\sim}J'^n, \tag{1.2.3}\\
 \Tor_i^A(J^n,A')&=0\quad\text{\uline{pour} }i\geq1. \tag{1.2.4}
\end{align}

a) On suppose d'abord que les $J^n/J^{n+1}$ sont projectifs sur $B=A/J$; comme $J$ est de type fini, ils sont localement libres de type fini sur $B$.
